% Глава 2: Обзор литературы / Chapter 2: Literature Review

\chapter{Обзор литературы}
\label{ch:literature}

\section{Наукометрические системы}
\label{sec:lit:scientometric}

[TODO: Обзор наукометрических систем: Web of Science, Scopus, ИСТИНА, Google Scholar]

\subsection{Международные системы}
\label{subsec:lit:international}

[TODO: Web of Science, Scopus, PubMed, arXiv]

\subsection{Национальные системы}
\label{subsec:lit:national}

[TODO: ИСТИНА (Россия), CERIF (Европа), Pure (Elsevier)]

\section{Проблема обработки многоязычных имен}
\label{sec:lit:multilingual}

[TODO: Обзор проблем обработки китайских, корейских, арабских имен]

\subsection{Китайские имена в научных публикациях}
\label{subsec:lit:chinese_names}

[TODO: Обзор исследований по порядку имени и фамилии, транслитерации, вариативности]

\subsection{Методы транслитерации и романизации}
\label{subsec:lit:transliteration}

[TODO: Pinyin, Wade-Giles, вариативность написания]

\section{Алгоритмы верификации данных}
\label{sec:lit:verification}

[TODO: Обзор методов верификации: rule-based, ML-based, hybrid approaches]

\subsection{Правиловые системы}
\label{subsec:lit:rule_based}

[TODO: Обзор rule-based approaches]

\subsection{Машинное обучение}
\label{subsec:lit:ml}

[TODO: Обзор ML подходов: supervised, unsupervised, transfer learning]

\subsection{Гибридные подходы}
\label{subsec:lit:hybrid}

[TODO: Комбинация правил и ML]

\section{Интерпретируемость систем принятия решений}
\label{sec:lit:interpretability}

[TODO: Обзор методов обеспечения интерпретируемости: LIME, SHAP, rule extraction]

\section{Защита персональных данных}
\label{sec:lit:privacy}

[TODO: Обзор методов защиты PII: hashing, anonymization, differential privacy]

\subsection{GDPR и ФЗ-152}
\label{subsec:lit:gdpr}

[TODO: Требования GDPR и Федерального закона 152-ФЗ]

\section{Выводы}
\label{sec:lit:conclusions}

[TODO: Выводы из обзора литературы, обоснование необходимости нового подхода]
